\section{GCC\EMDASH{}\EN{one more thing}\RU{ещё кое-что}\ES{una cosa más}}
\label{use_parts_of_C_strings}

\RU{Тот факт, что \IT{анонимная} Си-строка имеет тип}\EN{The fact that an \IT{anonymous} C-string has} 
\IT{const}\EN{ type} (\myref{string_is_const_char}), 
\RU{и тот факт, что выделенные в сегменте констант Си-строки гаратировано неизменяемые (immutable), 
ведет к интересному следствию}\EN{and the 
fact that C-strings allocated in constants segment are guaranteed to be immutable, has an interesting consequence}:\ES{El hecho de que una cadena de C \IT{anónima} sea de tipo \IT{const} (\myref{string_is_const_char}), y el hecho de que las cadenas de C asignadas en el segmento de constantes tengan garantizada su inmutabilidad, tiene una consecuencia interesante:}
\RU{компилятор может использовать определенную часть строки}\EN{the compiler may use a specific part of string}.\ES{el compilador puede usar una parte específica de la cadena.}

\RU{Вот простой пример}\EN{Let's try this example}\ES{Probemos este ejemplo}:

\begin{lstlisting}
#include <stdio.h>

int f1()
{
	printf ("world\n");
}

int f2()
{
	printf ("hello world\n");
}

int main()
{
	f1();
	f2();
}
\end{lstlisting}

\RU{Среднестатистический компилятор с \CCpp (включая MSVC) выделит место для двух строк, 
но вот что делает GCC 4.8.1}%
\EN{Common \CCpp{}-compilers (including MSVC) allocates two strings, 
but let's see what GCC 4.8.1 does}:\ES{Los compiladores comunes de \CCpp (incluido MSVC) asignan dos cadenas, pero veamos qué hace GCC 4.8.1}:

\begin{lstlisting}[caption=GCC 4.8.1 + \RU{листинг в }IDA\EN{ listing}\ES{listado en IDA}]
f1              proc near

s               = dword ptr -1Ch

                sub     esp, 1Ch
                mov     [esp+1Ch+s], offset s ; "world\n"
                call    _puts
                add     esp, 1Ch
                retn
f1              endp

f2              proc near

s               = dword ptr -1Ch

                sub     esp, 1Ch
                mov     [esp+1Ch+s], offset aHello ; "hello "
                call    _puts
                add     esp, 1Ch
                retn
f2              endp

aHello          db 'hello '
s               db 'world',0xa,0
\end{lstlisting}

\RU{Действительно, когда мы выводим строку}\EN{Indeed: when we print the \q{hello world} string}, 
\RU{эти два слова расположены в памяти впритык друг к другу и \puts, вызываясь из функции f2(), вообще не знает,
что эти строки разделены}\EN{these two words are positioned in memory adjacently and \puts called from f2() 
function is not aware this string is divided}.\ES{En efecto: cuando imprimimos la cadena \q{hello world}, estas dos palabras se colocan en memoria de forma adyacente, y \puts, al ser llamada desde la función f2(), no es consciente de que esta cadena está dividida.} \RU{Они и не разделены на самом деле, они разделены
только \q{виртуально}, в нашем листинге}\EN{It's not divided in fact, it's divided only \q{virtually}, in this
listing}.\ES{De hecho, no está dividida; está dividida solo \q{virtualmente} en este listado.}

\RU{Когда}\EN{When}\ES{Cuando} \puts \RU{вызывается из f1(), он использует строку}\EN{is called from f1(), it uses}\ES{es llamada desde f1(), utiliza la cadena} 
\q{world} \RU{плюс нулевой байт}\EN{string plus zero byte}.\ES{ más el byte cero.} \puts \RU{не знает, что там ещё есть какая-то строка
перед этой}\EN{is not aware there is something before this string}!\ES{¡\puts no es consciente de que hay algo antes de esta cadena!}

\RU{Этот трюк часто используется (по крайней мере в GCC) и может сэкономить немного памяти.}
\EN{This clever trick is often used by at least GCC and can save some memory.}\ES{Este ingenioso truco suele ser utilizado al menos por GCC y puede ahorrar algo de memoria.}
